\documentclass[11pt,a4paper]{article}

% ---------------------------------------------------------
% PREAMBLE (stable, warning-free, Overleaf-safe)
% ---------------------------------------------------------
\usepackage[utf8]{inputenc}
\usepackage[T1]{fontenc}
\usepackage{lmodern}
\usepackage{amsmath, amssymb, amsfonts}
\usepackage{bm}
\usepackage{geometry}
\usepackage{graphicx}
\usepackage{hyperref}
\usepackage{cite}
\usepackage{authblk}
\usepackage{physics}
\usepackage{mathtools}

\geometry{margin=2.4cm}

\title{\textbf{Companion LXXII: Kernel PCA of Persistence Diagrams in the Relaxation-Driven Cyclic Cosmology}}
\author{Michael Lehmann \\ \texttt{mi.lehmann@gmx.de}}
\date{April 2026}

\begin{document}
\maketitle

\begin{abstract}
Kernel principal component analysis (kernel PCA) provides a spectral decomposition 
of persistence diagrams embedded in a reproducing kernel Hilbert space (RKHS). In 
weak-lensing analyses, this decomposition reveals the dominant modes of 
topological variation across connected components, loops, and void-like structures. 
In the standard cosmological model, the kernel spectrum reflects the nonlinear 
matter distribution and the projection of large-scale structure. In the 
Relaxation-Driven Cyclic Cosmology (RDCC), the scale-dependent evolution of the 
gravitational potential modifies the kernel eigenvalue spectrum in a characteristic 
way. This Companion develops a continuous description of kernel PCA in RDCC, 
deriving how the relaxation scale influences spectral topology, principal 
components, and the observational signatures in weak-lensing surveys. The 
resulting framework provides a distinctive probe of the RDCC gravitational field 
through spectral decomposition of linearized topology.
\end{abstract}
% ---------------------------------------------------------
\section{Introduction}

Persistence diagrams encode the birth and death of topological features in 
weak-lensing convergence fields. Hilbert-space embeddings provide a linearized 
representation of these diagrams, enabling kernel methods and inner-product 
computations. Kernel PCA extends this framework by decomposing the embedded 
topology into orthogonal spectral modes.

In the standard cosmological model, the kernel eigenvalue spectrum reflects the 
nonlinear matter distribution and the projection of large-scale structure. In 
RDCC, the gravitational potential evolves differently. The relaxation mechanism 
suppresses the decay of small-scale modes and enhances the coherence of large-scale 
modes. This modifies the kernel spectrum in a scale-dependent manner, producing a 
distinctive signature in kernel PCA analyses.
% ---------------------------------------------------------
\section{Kernel PCA of Persistence Diagrams}

Given persistence diagrams $D_{1}, \dots, D_{N}$ embedded into an RKHS via 
$\Phi(D_{i})$, the kernel matrix is

\begin{equation}
    K_{ij}
    = \langle \Phi(D_{i}), \Phi(D_{j}) \rangle_{\mathcal{H}}.
\end{equation}

Kernel PCA solves the eigenvalue problem

\begin{equation}
    K \bm{v}_{k}
    = \lambda_{k} \bm{v}_{k},
\end{equation}

where $\lambda_{k}$ are the kernel eigenvalues and $\bm{v}_{k}$ the corresponding 
eigenvectors.

The $k$-th principal component is

\begin{equation}
    \mathrm{PC}_{k}(D_{i})
    = \sum_{j=1}^{N} v_{k,j} K_{ij}.
\end{equation}
% ---------------------------------------------------------
\section{RDCC Modification of Kernel Structure}

In RDCC, the convergence evolves as
\begin{equation}
    \kappa_{\mathrm{RDCC}}(\ell)
    = \kappa(\ell)
      \left[
        1 - \chi_{\kappa}
        \left(\frac{\ell}{\ell_{\mathrm{rel}}}\right)^{\alpha}
      \right].
\end{equation}

This modifies the birth and death thresholds of topological features, and hence the 
kernel:

\begin{equation}
    K^{\mathrm{RDCC}}_{ij}
    = K_{ij}
      \left[
        1 - \chi_{K}
        \left(\frac{\ell}{\ell_{\mathrm{rel}}}\right)^{\alpha}
      \right].
\end{equation}

The RDCC eigenvalue problem becomes

\begin{equation}
    K^{\mathrm{RDCC}} \bm{v}^{\mathrm{RDCC}}_{k}
    = \lambda^{\mathrm{RDCC}}_{k} \bm{v}^{\mathrm{RDCC}}_{k}.
\end{equation}
% ---------------------------------------------------------
\section{Spectral Signatures in RDCC}

To first order in the relaxation parameter,

\begin{equation}
    \lambda^{\mathrm{RDCC}}_{k}
    = \lambda_{k}
      \left[
        1 - \chi_{\lambda,k}
        \left(\frac{\ell}{\ell_{\mathrm{rel}}}\right)^{\alpha}
      \right].
\end{equation}

RDCC therefore induces:

\begin{itemize}
    \item enhanced large-scale eigenvalues (coherent modes),
    \item suppressed small-scale eigenvalues (relaxation damping),
    \item a characteristic turnover near $\ell_{\mathrm{rel}}$,
    \item modified ordering of principal components,
    \item altered projection of persistence lifetimes onto spectral modes.
\end{itemize}
% ---------------------------------------------------------
\section{Principal Components in RDCC}

The RDCC principal components are

\begin{equation}
    \mathrm{PC}^{\mathrm{RDCC}}_{k}(D_{i})
    = \sum_{j=1}^{N}
      v^{\mathrm{RDCC}}_{k,j}
      K^{\mathrm{RDCC}}_{ij}.
\end{equation}

These components encode the dominant RDCC-modified topological modes.
% ---------------------------------------------------------
\section{Conclusion}

Kernel PCA of persistence diagrams in the Relaxation-Driven Cyclic Cosmology 
reveals the spectral structure of linearized topology. The relaxation mechanism 
modifies the kernel eigenvalue spectrum in a scale-dependent manner, producing 
distinctive signatures in weak-lensing surveys. These effects provide a direct 
observational window into the relaxation scale and the temporal structure of the 
RDCC gravitational field. The present Companion establishes the theoretical 
foundation for future studies of spectral topology, kernel methods, and 
statistical tests in RDCC.

\bibliographystyle{apsrev4-2}
\nocite{*}
\bibliography{references}

\end{document}
